\documentclass[11pt,a4paper]{article}

\usepackage[utf8]{inputenc}
\usepackage[T1]{fontenc}
\usepackage{lmodern}
\usepackage[margin=1in]{geometry}
\usepackage{microtype}
\usepackage{booktabs}
\usepackage{tabularx}
\usepackage{xcolor}
\usepackage{enumitem}
\usepackage{titlesec}
\usepackage{abstract}
\usepackage[hidelinks]{hyperref}
\usepackage{caption}

\definecolor{navy}{HTML}{1F3A5F}
\definecolor{accent}{HTML}{B5791F}

\hypersetup{colorlinks=true, linkcolor=navy, citecolor=navy, urlcolor=navy}

\titleformat{\section}{\normalfont\large\bfseries\color{navy}}{\thesection}{0.6em}{}
\titleformat{\subsection}{\normalfont\normalsize\bfseries\color{navy}}{\thesubsection}{0.6em}{}
\setlist{itemsep=2pt, topsep=3pt, parsep=0pt}
\newcolumntype{Y}{>{\raggedright\arraybackslash}X}
\renewcommand{\arraystretch}{1.2}

\title{\vspace{-1.0cm}\textbf{\color{navy} Do Skills Make ML-Engineering Agents\\[2pt] Measurably Better?}\\[6pt]
\large A System for Building, Discovering, and Evaluating Agent Skills}

\author{Kuntal Kokate\\[2pt]
\normalsize PengTao Lab, University of California, San Diego\\
\normalsize \texttt{kukokate@ucsd.edu}}

\date{June 2026}

\begin{document}
\maketitle

\begin{abstract}
\noindent
Agent Skills---structured Markdown playbooks injected into a language-model agent's
context---are increasingly used to specialise general agents for narrow domains.
Yet there is little evidence that a skill changes what an agent \emph{builds},
as opposed to what it \emph{says}. This report presents an end-to-end system that
(i)~\textbf{builds} a progressive-disclosure skill from a single package URL,
(ii)~\textbf{discovers} and safely installs existing community skills, and
(iii)~\textbf{evaluates} whether a skill makes a machine-learning-engineering (MLE)
agent measurably better. The intended end state is a \emph{library} of builder-made
skills handed to the agent, from which it selects a few of its own choosing while
solving a task; the evaluation framework is built to test precisely that. The
evaluation is a two-stage pipeline: a fast, CI-style
local test that asks whether a skill \emph{fires} and surfaces the right facts, and
a paired with-skill / without-skill A/B framework that runs a real MLE search agent
(MLEvolve) on parameter-efficient fine-tuning (PEFT) tasks and grades the outcome
with an \emph{independent held-out grader}. We describe the architecture, the
progressive-disclosure injection mechanism, the trustworthy-metric design, and the
Kubernetes runtime, and we report preliminary SAMSum results from the held-out grader:
in a four-cell $2\times2$ run, two cells scored validly (ROUGE-L $0.29$ with-skill and
$0.23$ without-skill, on different seeds) while an output-contract bug invalidated the
other two, so no paired lift can yet be claimed. We close with the honest
limitations of these spike runs and the design of the full sweep.
\end{abstract}

\section{Introduction}

A general-purpose coding agent can be specialised by injecting an \emph{Agent Skill}:
a Markdown document describing when and how to use a particular library or technique,
following a \emph{progressive-disclosure} convention in which a short description is
always loaded, a body is loaded on trigger, and reference files are loaded only as
needed. Skills are cheap to write and easy to share, but their effect is usually
asserted rather than measured. The central question of this project is deliberately
operational:

\begin{quote}
\itshape Does an MLE agent build a measurably better pipeline \emph{with} a skill than
\emph{without} it---and, if so, \emph{where} in the pipeline does the help land?
\end{quote}

Answering this requires three capabilities, which map onto the three loosely-coupled
bodies of work in this project: a way to \textbf{build} good skills, a way to
\textbf{discover} existing ones, and---the research focus---a rigorous way to
\textbf{evaluate} their effect on a downstream agent.

\paragraph{The intended end state.} The three capabilities are not independent
deliverables but stages of one loop. The builder produces many skills into a shared
\emph{library}; at task time the \emph{entire} library---not a single hand-picked
skill---is made available to the MLE agent, which selects a few skills of its own
choosing as it works, loading only what each step needs. The evaluation framework
exists to test exactly this arrangement: whether a builder-made library with
agent-driven selection makes the agent measurably better, and whether the agent
selects the \emph{right} skills.

The remainder of this report describes each capability in turn, with the bulk of the
discussion devoted to the evaluation framework.

\section{System Overview: Three Agents}

The build and discovery capabilities are implemented as two OpenClaw agents; a third
agent supports the local evaluation stage. Table~\ref{tab:agents} summarises them.

\begin{table}[h]
\centering
\caption{The three agents and their roles.}
\label{tab:agents}
\small
\begin{tabularx}{\textwidth}{l Y l}
\toprule
\textbf{Agent} & \textbf{Role} & \textbf{Status} \\
\midrule
\texttt{ai-skill-builder} & Synthesises a \texttt{SKILL.md} (plus references and scripts) from one documentation URL. & Phase 1.5, v1.4.0 \\
\texttt{ai-skill-scout}   & Searches GitHub for existing skills, ranks by trust, and installs them after a security scan. & Phase 1, complete \\
\texttt{skill-tester}     & Drives the Stage-1 local evaluation; an MCP sidecar captures every prompt for measurement. & In production \\
\bottomrule
\end{tabularx}
\end{table}

\subsection{Skill-Builder}

The builder turns one URL into one skill. Its pipeline is:
\textsc{resolve} the owner/repo from the URL; \textsc{fetch} documentation, README, and
examples (optionally closed bug issues and Stack~Exchange Q\&A); \textsc{extract}
hardware hints with a regex side-channel; \textsc{scan} community text for
indirect-prompt-injection (IPI); \textsc{plan} the file structure with one LLM call;
\textsc{synthesise} the body, reference files, templates, and evals in parallel;
\textsc{validate} (a 60-pattern security scanner across seven threat categories, a
line cap, a shell-command-fabrication check, and \texttt{py\_compile} on templates);
run a \textsc{triggering eval} that judges the description against decoys and rewrites
it once if it loses; \textsc{assemble} the YAML frontmatter \emph{deterministically in
Python}; and \textsc{write} the bundle to the workspace.

Two design choices matter. First, \textbf{the model never writes the frontmatter}: the
\texttt{name}, install steps, runtime MCP declarations, and a content hash for
provenance are assembled by code, removing a class of hallucination. Second,
\textbf{forum text is excluded by design}; Stack~Exchange (licensed, vote-deduplicated,
attribution-preserving) and closed issues are the community substrate, because
unmoderated forums carry disproportionate prompt-injection risk. On the canonical
\texttt{peft-tuning} skill the builder's triggering test scored an F1 of $1.000$
(10 should-fire and 10 near-miss prompts, 60 judge calls).

\subsection{Skill-Scout}

The scout's source is GitHub only. It expands the user's query into synonyms via an
LLM, merges \texttt{gh search} results, fetches each candidate's full \texttt{SKILL.md},
and scores \emph{trust} (known orgs / star count), \emph{completeness} (0--10), and a
\emph{classification} (\textsc{adopt} / \textsc{extend} / \textsc{gap}). Candidates are
ranked by \texttt{(trust, class, completeness, stars)} so a high-trust skill that needs
extending outranks a polished one of unknown origin. Crucially, nothing reaches the
workspace until it has been quarantined to \texttt{/tmp} and cleared the \emph{same}
60-pattern scanner the builder uses---one source of truth for what is dangerous---and
the user has approved the candidate.

\subsection{Skill-Tester}

The tester agent supports Stage~1 of the evaluation (Section~\ref{sec:stage1}). An MCP
sidecar captures native tool calls, a bash side-channel log, and narration so that
triggering and functional lift can be measured against a baseline agent, and reply text
is preserved for offline re-grading.

\section{Evaluation Methodology}

The evaluation is a two-stage pipeline. Stage~1 asks a cheap question in isolation;
Stage~2 asks the expensive question inside a full agent run. The division lets us reject
weak skills before spending GPU hours.

\subsection{Stage 1: Local Skill Evaluation}
\label{sec:stage1}

Stage~1 runs on every build (CI-style, $\approx 5$ minutes, $\approx \$1$) and follows
the Anthropic \texttt{skill-creator} protocol. It reports four metrics: a
\textbf{triggering F1} over a 20-prompt 60/40 should-fire / near-miss split (target
$\geq 0.85$); a \textbf{functional pass-rate} on prompts with deterministic
\texttt{must\_contain} assertions, run with and without the skill; a \textbf{lift} (the
difference between the two); a \textbf{citation rate} (does the reply name a reference
file; target $\geq 80\%$); and a \textbf{token-cost ratio} (target $\leq 2\times$).
A known limitation is \emph{saturation}: on canonical questions the base model already
knows the answer, so deterministic checks pass with or without the skill, and the lift
appears only on harder, version-specific questions---the very effect Stage~2 is designed
to surface.

\subsection{Stage 2: The MLEvolve A/B Framework}

Stage~2 is the core contribution. It runs a real MLE agent on a freeform task twice---once
with the skill injected and once without---holding everything else fixed, and measures the
difference.

\paragraph{Agent.} The agent is \textbf{MLEvolve-generic}, a Monte-Carlo graph search
that generates, executes, and evolves candidate code nodes (a fork of AutoMLGen's search,
pinned to a fixed commit), driven by \texttt{deepseek/deepseek-v4-pro} via OpenRouter. It
replaced an earlier agent because its \emph{subprocess-per-node} execution
avoids a fork-after-CUDA failure mode that crashed the GPU pod.

\paragraph{Experimental design.} Every controllable variable is held constant across the
two cells: the same backbone (\texttt{Qwen/Qwen2.5-3B-Instruct}), the same task, the same
random seed, and the same skill \emph{library}. The \emph{only} difference is whether that
library is made available to the agent; when it is, the agent selects from it, and an empty
library reproduces the baseline exactly. The current single-seed spike runs held one skill
(\texttt{peft-tuning}) fixed---a ``skill vs.\ no skill'' contrast---whereas the designed end
state is a multi-skill library from which the agent chooses, which additionally exposes a
\emph{selection-quality} signal (described below). The task pool is three contract-only PEFT
fine-tuning tasks (Table~\ref{tab:tasks}); the full sweep is $3$ seeds $\times\ 3$ tasks
$\times\ 2$ cells $= 18$ trajectories, with a 95\% confidence interval estimated from
paired-seed variance.

\begin{table}[h]
\centering
\caption{Stage-2 task pool.}
\label{tab:tasks}
\small
\begin{tabularx}{\textwidth}{l Y l}
\toprule
\textbf{Task} & \textbf{Domain} & \textbf{Metric} \\
\midrule
\texttt{samsum} & Dialogue summarisation (SFT) & ROUGE-L F1 \\
\texttt{gsm8k}  & Math reasoning (SFT)         & Exact-match accuracy \\
\texttt{boolq}  & Yes/no question answering (SFT) & Accuracy (test labels withheld) \\
\bottomrule
\end{tabularx}
\end{table}

\paragraph{Skill injection by progressive disclosure.} MLEvolve performs single-shot code
generation and cannot open a file to read a skill on its own. The injector stands in for
that, mapping Anthropic's Discovery $\rightarrow$ Activation $\rightarrow$ Execution model
onto the single seam (\texttt{get\_impl\_guideline\_from\_agent}) that every code-generation
agent calls, via a \texttt{sys.meta\_path} import hook that patches all four agents
(draft, improve, debug, evolution):
\begin{itemize}
\item \textbf{Tier 0 --- Discovery (every node).} A catalogue---each skill's name, one-line
description, and reference filenames ($\approx 150$ tokens for the whole library)---is
appended to the implementation guideline, so the agent is always aware of what exists.
\item \textbf{Tier 1 --- Activation (per node).} A temperature-0 selector runs once per node,
keyed to the current stage, and decides which skill(s) this step actually needs.
\item \textbf{Tier 2 --- Execution.} The same selector chooses which \texttt{references/*.md}
files to splice in---never the whole library, never every body into every node.
\end{itemize}
When the skill library is empty (the without-skill cell), there is no catalogue and no
selector call, so the guideline passes through byte-for-byte unchanged---guaranteeing the
two cells differ only in the treatment.

\paragraph{Trustworthy metric.} An agent that reports its own score will learn to report a
good one. Therefore the headline number never comes from the agent. After the run exits, an
\emph{independent held-out grader} (\texttt{mleval.grader}) recomputes the task metric from
the agent's preserved \texttt{submission.csv} against references the agent was never shown.
MLEvolve's own self-reported validation score is retained only as the tree-search signal and
a drift diagnostic; in one run a self-reported $0.81$ graded out as invalid, which is exactly
the gaming this design prevents.

\paragraph{Measurement layers.} Beyond the L1 outcome (held-out grader $+$ paired
Lift), an \textbf{L2} layer attributes \emph{where} the skill helped using a $6\times16$
pipeline-stage taxonomy and three co-location-proof metrics (clean-reach, rework,
failure-modes) derived from an AST classifier over the agent's code; an \textbf{L3} layer
records wall-clock, token, and dollar cost so a skill that helps but slows the agent is
not mistaken for a free win. Finally, because the injector logs every
\texttt{select\_skills} call, a \textbf{selection-quality} layer measures whether the agent
chose the right skills---the precision and recall of its per-task selections against a
hand-labelled oracle of which library skills are genuinely relevant to each task. This
layer is meaningful only once the library holds more than one skill, and becomes the
headline as the project moves from the single-skill contrast to the full library.

\section{Implementation and Runtime}

The framework runs on the UCSD Nautilus NRP Kubernetes cluster: one container image per
agent plugin, and one Pod per \texttt{(task $\times$ cell $\times$ seed)} trajectory. The
MLEvolve agent is adapted at runtime by a small set of import-time monkey-patches (a
``sidecar''), kept to the minimum needed for a fair, reproducible A/B:
\begin{itemize}
\item \texttt{seed} --- pins \texttt{random} / \texttt{numpy} / \texttt{torch} RNG from the
trajectory's seed, enabling paired-seed comparison;
\item \texttt{prompt\_logger} --- captures every prompt, response, and token count;
\item \texttt{token\_budget} --- raises the default output cap $16{,}384 \rightarrow 32{,}768$
to stop a mid-output truncation that corrupted generated code;
\item \texttt{skill\_injector} --- the only patch that carries the treatment (Tier-0/1/2 above).
\end{itemize}
Separately, a build-time patch neutralises an upstream ``Kaggle grandmaster'' persona that
was driving the agent off-task onto the wrong dataset. The container entrypoint handles
signals so that the analyzer chain always runs and the trajectory's outputs always land on
the shared PVC, even on a deadline kill.

\section{Preliminary Results}

The harness has been validated end-to-end across a series of SAMSum spike runs; the most
complete is \texttt{spike-018}, which ran the full $2\times2$ (with/without skill $\times$
seeds 0--1) under the independent held-out grader. Of the four cells, only two produced a
submission the grader could score---an output-contract bug (a submission with the wrong
columns) invalidated one cell in \emph{each} arm---so there is no seed-matched pair and
\textbf{no Lift can yet be computed}. The two valid held-out ROUGE-L scores are $0.288$
(with-skill, seed~1) and $0.227$ (without-skill, seed~0); coming from different seeds, they
are indicative only. The grader also rejected two inflated self-reported scores ($0.94$ and
$0.81$) as invalid---precisely the off-task drift it was built to catch (an earlier
\texttt{spike-012} smoke had reported $0.4331$ on a self-reported metric whose control had
drifted to a different task). On \emph{process} quality the with-skill arm is consistently
stronger: a self-correction rate of $\approx\!0.6$ versus $\approx\!0.15$, markedly less
redundant work (3--4 versus 9--12 redundant nodes), and a significant pipeline-stage
redistribution ($\chi^2 = 65.1$, $p \approx 3\times10^{-9}$).

\paragraph{gsm8k.} The \texttt{gsm8k} task was rewritten to an MLE-Bench-aligned held-out
re-split (labels withheld on the test set; a carved validation slice for the agent's own
signal). Run \texttt{spike-023} then validated the held-out pipeline end-to-end: a
without-skill node produced a valid 1{,}319-row submission scoring \textbf{0.61 exact-match}
($809/1{,}319$) under the held-out grader (graded from the preserved submission; the
graceful-kill path did not auto-persist the score file). Two findings shaped the next
iteration. First, a skill-router regression left the treatment arm effectively empty: an
$\approx$3\,KB harness-rules header pushed the task's model, method, and data past the
selector's 1{,}500-character cap, so the per-node selector returned no skills and no skill
body ever reached a code-generation prompt; this was fixed with an explicit end-of-rules
sentinel and a larger cap. Second, the binding constraint is \emph{generation} throughput,
not training: every node trains in $\approx$10--15 minutes, and timeouts occur in the
1{,}319-example decode---one node had to ratchet its batch size $4\!\rightarrow\!32$ and
\texttt{max\_new\_tokens} $512\!\rightarrow\!200$ to finish within the 52-minute window.
Because our harness imposes a hard 60-minute per-execution cap and marks any unfinished
execution buggy (no partial credit), slowness is punished categorically; the planned response
is to raise the cap rather than teach batching tricks in the instruction, which would erase
the very signal the skill is meant to supply. No clean paired gsm8k A/B has completed yet.

These numbers should be read as \emph{spike findings, not the verdict}. They come from a
debugging-era sweep rather than the full paired 18-trajectory run, and the decisive failures
hit \emph{both} cells---the output-contract bug above (a wrong id column), fork-after-CUDA
segmentation faults, and a period in which the skill router was blind to the task because the
harness rules pushed the task description past the selector's character budget. Each was
diagnosed and fixed; the harness now grades, retries, and harvests trajectories cleanly, so a
clean paired Lift awaits a rerun.

\section{Issues Surfaced in the Skill Builder}

Running the skills inside a live agent---rather than only triggering them in isolation---exposed
limitations in the builder that a triggering test cannot see. The clearest case is the
\texttt{vllm-inference} skill on \texttt{gsm8k}: it passed the builder's triggering eval
(F1 $\geq 0.85$; the live selector picked it on 10 of 11 nodes) yet made no positive
behavioural difference. Three structural gaps explain the distance between ``fires'' and
``helps'':
\begin{itemize}
\item \textbf{Wrong tool, correctly avoided.} The task prescribes Hugging~Face
\texttt{AutoModelForCausalLM} with batched generation; vLLM is never intended. The agent used
Hugging~Face in \emph{both} arms, so the earlier reading that ``the agent ignored the vLLM
skill'' was a mis-diagnosis, not a skill defect.
\item \textbf{Genre self-exclusion.} The skill's \emph{NOT-for} section directs supervised and
LoRA fine-tuning to \texttt{transformers}; because \texttt{gsm8k} is a LoRA task, the skill
scoped \emph{itself} out and steered the agent away---counterproductive rather than merely
inert---despite its own decision tree listing offline evaluation as in-scope.
\item \textbf{The decisive failure was functional, not selection.} Both arms ran Hugging~Face
regardless of what was selected; the runs failed on generation \emph{hygiene}---\texttt{use\_cache}
disabled by gradient checkpointing, loose \texttt{max\_new\_tokens}, and validation generation
every epoch---which the co-resident \texttt{peft-tuning} skill (an older \texttt{1.3.0} build)
documents as API surface, not as workflow choreography.
\end{itemize}

The triggering eval missed this because its judge competes the new skill against five canned
generic decoys rather than its real co-resident siblings, so the disambiguation that
mattered---fine-tune-and-evaluate should lead with \texttt{peft-tuning}, not
\texttt{vllm-inference}---was never tested, and because the judge scores short user messages
whereas the live selector reads a multi-paragraph task instruction. Two creator-level fixes
follow: (1)~a \emph{library-aware} triggering eval that judges a skill against its actual
siblings, including hard fine-tune-versus-serve cases; and (2)~a \emph{functional}
(choreography) eval that rewards a skill for encoding operational gotchas, not merely for
owning the right keywords. Tightening descriptions alone would not have caught this.

\section{Discussion and Limitations}

The clearest lesson is methodological: a credible skill-effect number requires a grader the
agent cannot reach, and a treatment whose only difference from the baseline is the skill
itself. A second lesson, consistent with the literature, is that the skill's
\textbf{description} is its single highest-leverage component---wording alone swings
invocation by more than an order of magnitude with capability held fixed---so a skill that
never fires zeroes out every other quality it has. The \texttt{vllm-inference} case (preceding
section) sharpens this in the other direction: discovery is \emph{necessary but not
sufficient}---a skill can be selected on nearly every node and still fail to help---so a
credible evaluation must reach past triggering into functional behaviour. Limitations remain:
the AST stage
classifier is a rule-table MVP pending a call-graph upgrade; the LLM judge for subjective
assertions is off by default; and the full paired sweep, with confidence intervals, has not
yet been run.

\section{Conclusion and Future Work}

This project delivers a closed loop for Agent Skills---build, discover, and evaluate---with
the evaluation framework as the research contribution. The immediate next step is the full
paired A/B across \texttt{samsum}, \texttt{gsm8k}, and \texttt{boolq}, reporting mean Lift
with a 95\% confidence interval, the L2 stage-level attribution, and---once the library
holds several builder-made skills---the agent's skill-selection precision and recall against
a per-task oracle. The promotion rule is
deliberately conservative: a skill is shipped only when it lifts \emph{both} task outcome and
recipe fidelity, never one at the expense of the other. In parallel, the builder's own
evaluation will be hardened along the two axes the \texttt{vllm-inference} case exposed---a
library-aware triggering eval and a functional choreography check---so that a skill which
\emph{fires} reliably also \emph{helps}.

\end{document}
